% ========================================================================
% In-house reference paper. Kept in this repo as the baseline against
% which the IMD-augmented extension is benchmarked.
%
%   Authors : G. Pallares, M. Foucras, R. Fossé, I. Pabois (CESI LINEACT)
%   Venue   : SCITEPRESS conference format (draft 2026-05)
%   Task    : 15-min station-level bike availability forecast
%   Data    : 15 French cities, 1-min station_status, Sep-Nov 2024
%             primary case study: "Place Saint-Pierre, Toulouse"
%   Features: AB (available bikes), MM (min since midnight),
%             WD (weekday), RH (rel. humidity), T (temperature)
%   Models  : LSTM, MLP, ARIMA(2,0,2), XGBoost (200 trees, depth 6)
%   Best    : LSTM, MAE = 1.85 (~8.5% of 22-slot capacity) at 15 min,
%             AB+MM+WD only, 2-month training, 1-month test
%   Key finding: weather degrades performance ; recent usage dominates
%
% Extension plan (paper_demand/related/IMPROVEMENT_PLAN.md) :
%   1. Add IMD-4 station features (M, I, T, D + 9 socio-econ)
%   2. Add Transformer / TFT / Chronos baselines
%   3. Extend from 1 station to all 15 cities × all stations
%   4. Add cross-city transfer test
%   5. Add bootstrap CIs on MAE / RMSE
%   6. Re-validate the "weather doesn't help" claim with rainfall data
% ========================================================================

\documentclass[a4paper,twoside]{article}
\usepackage{epsfig}
\usepackage{subcaption}
\usepackage{calc}
\usepackage{amssymb,amstext,amsmath,amsthm}
\usepackage{multicol}
\usepackage{pslatex}
\usepackage{booktabs}
\usepackage{algorithm2e}
\usepackage[bottom]{footmisc}
\usepackage{url}
\usepackage{xcolor}
\usepackage{comment}
% \usepackage{SCITEPRESS}  % requires venue-side .sty

\newcommand{\ie}{\emph{i.e.}}
\newcommand{\eg}{\emph{e.g.}}
\newcommand{\etc}{etc{.\@}}

\begin{document}
\title{Forecasting Short-Term Availability of Bikes at Sharing Stations:
A Deep-Learning Approach}

\author{Ga\"el Pallares, Myriam Foucras, Rohan Foss\'e, Ilan Pabois\\
CESI LINEACT, France}

\date{2026}
\maketitle

\begin{abstract}
The growing popularity of bike-sharing systems represents a significant
urban mobility trend. In this paper, the short-term bike availability
forecasting at the station level is addressed using real-time open data
from public bike station feeds across many French cities. We show that
deep learning methods (LSTM, MLP) achieve a mean forecasting error
lower than 10\,\% relative to station capacity even for limited training
data volumes and number of features. Results are computed in reduced
time, opening the way for intelligent fleet management.
\end{abstract}

\paragraph{Data.} 15 French cities, station\_status at 1-min resolution,
plus hourly weather (T, RH) and hourly air-pollution (NO$_2$, O$_3$,
PM10, PM2.5, SO$_2$, NO$_x$, C$_6$H$_6$, CO, NO). The principal case
study is the station ``Place Saint-Pierre'' in Toulouse, 22-slot
capacity.

\paragraph{Features (5).}
AB = available bikes ;
MM = minutes since midnight ;
WD = weekday ;
RH = relative humidity ;
T = temperature.

\paragraph{Models.}
Two-layer MLP (100/100, ReLU). Two-layer LSTM (100/100 memory units,
dropout 0.2). Both: Adam, lr=1e-3, MSE loss, batch 32, $\sim 300$
epochs with early stopping. Plus ARIMA$(2,0,2)$ rolling forecast and
XGBoost (200 trees, depth 6, lr 0.1).

\paragraph{Best result.}
Place Saint-Pierre, LSTM, AB+MM+WD only, 2-month training,
1-month test : MAE = 1.85 bikes ($\approx 8.5\,\%$ of capacity) at
15-min horizon ; MAE = 2.65 at 60-min horizon. Adding T or RH
\emph{degrades} performance. ARIMA and XGBoost converge to similar
MAE on long training windows ; XGBoost drops markedly below 15-day
training.

\paragraph{Key findings.}
\begin{itemize}
\item Recent autoregressive state dominates : AB+MM+WD is enough.
\item Slow exogenous covariates (T, RH) are redundant at $<1$\,h horizon.
\item Training beyond 1 month does not improve MAE (concept drift).
\item LSTM beats MLP at long training ; MLP is competitive at short
training and is faster.
\item The model predicts \emph{observed} availability, not unmet
demand : empty/full stations censor true user need.
\end{itemize}

\paragraph{Open extensions identified in the discussion.}
\begin{itemize}
\item Add precipitation and wind to revisit the weather claim.
\item Add SHAP / attention for interpretability.
\item Generalise across all 15 cities (currently single-station focus).
\item Compare against modern spatio-temporal architectures.
\end{itemize}

% End of summary. The full original draft is the working copy held by
% G. Pallares (gpallares@cesi.fr) ; this archival summary is the
% baseline reference for the IMD-augmented extension scoped in
% IMPROVEMENT_PLAN.md.

\end{document}
